\section{Finite-field counts and candidate Frobenius certificates}\label{sec:frobenius}

For the selected primes $p=5,11,13$, direct modular gcd checks leave only the
expected affine node, keep $Q_6$ squarefree and disjoint from that node, and
preserve the displayed leading infinity charts.  These checks support the
expected normalization pattern at the three reductions.  They do not
construct a smooth proper simultaneous normalization over $\mathbb Z_p$, prove that
normalization commutes with reduction, or classify all bad primes.

Let $A_{p,r}$ count affine solutions of $P=0$ over $\F_{p^r}$.  Applying the
characteristic-zero branch ledger formally adds seven rational branches at
infinity.  At the rational plane node, the two tangent directions are
rational exactly when $-7$ is a square in $\F_{p^r}$.  We therefore define
the explicitly reproducible \emph{branch-corrected count}
\begin{equation}\label{eq:countcorrection}
 \widehat N_{p,r}=A_{p,r}+7+\epsilon_{p,r},\qquad
 \epsilon_{p,r}=\begin{cases}+1,&-7\text{ is a square},\\-1,&\text{otherwise}.
 \end{cases}
\end{equation}
A simultaneous-normalization proof would promote $\widehat N_{p,r}$ to the
point count of the good reduction of $\Ccurve$; absent that proof, the hat is
part of the claim boundary.

\begin{table}[ht]
\centering
\caption{Exact affine and branch-corrected counts.}
\begin{tabular}{c@{\qquad}c@{\qquad}c}
\toprule
$p$ & $(A_{p,1},A_{p,2},A_{p,3})$ & $(\widehat N_{p,1},\widehat N_{p,2},\widehat N_{p,3})$\\
\midrule
5  & $(3,31,141)$    & $(9,39,147)$\\
11 & $(11,159,1163)$ & $(19,167,1171)$\\
13 & $(10,234,2125)$ & $(16,242,2131)$\\
\bottomrule
\end{tabular}
\end{table}

Set $\widehat S_r=p^r+1-\widehat N_{p,r}$.  Newton identities and the
genus-three reciprocal completion define the candidate numerator
\begin{equation}\label{eq:localnumerator}
\begin{split}
\widehat L_p(T)={}&1-\widehat S_1T+
\frac{\widehat S_1^2-\widehat S_2}{2}T^2
-\frac{\widehat S_1^3-3\widehat S_1\widehat S_2+2\widehat S_3}{6}T^3\\
&+p\frac{\widehat S_1^2-\widehat S_2}{2}T^4
-p^2\widehat S_1T^5+p^3T^6.
\end{split}
\end{equation}
The three candidates are
\begin{align*}
\widehat L_5(T)&=1+3T+11T^2+31T^3+55T^4+75T^5+125T^6,\\
\widehat L_{11}(T)&=1+7T+47T^2+161T^3+517T^4+847T^5+1331T^6,\\
\widehat L_{13}(T)&=1+2T+38T^2+51T^3+494T^4+338T^5+2197T^6.
\end{align*}
Each is irreducible over $\Q$, satisfies
$[T^{6-i}]\widehat L_p=p^{3-i}[T^i]\widehat L_p$, and has
reciprocal-root modulus residual below $2.2\cdot10^{-15}$ from $p^{-1/2}$.
This circle check is a regression property of the candidates, not a proof of
good reduction and not evidence for the Riemann hypothesis.

The producer uses the Python \texttt{galois} package.  The checker instead
constructs $\F_{p^r}$ as an explicit polynomial quotient, builds its own
addition and multiplication tables, and imports no producer code.  It repeats
all nine frozen affine counts and every branch correction.  The recurrence
from Eq.~\eqref{eq:localnumerator} predicts $\widehat N_{5,4}=547$; a sealed
direct count gives $A_{5,4}=539$ and Eq.~\eqref{eq:countcorrection} gives
exactly 547.
